% -----------------------------------------------------------------------
% CAPÍTULO 6: EXPERIMENTOS Y PRUEBAS
% -----------------------------------------------------------------------
\chapter{Experimentos y Pruebas}

Este capítulo describe el proceso de validación y verificación llevado a cabo para asegurar la calidad, estabilidad y correcto funcionamiento de la plataforma AIForge. Se detallan las pruebas realizadas tanto en la interfaz de usuario como en los servicios de backend y el modelo de inteligencia artificial.

% -----------------------------------------------------------------------
\section{Pruebas realizadas preproducción}
% -----------------------------------------------------------------------

Esta sección agrupa las pruebas realizadas antes del despliegue en producción: metodología y pruebas críticas del frontend, validación de la API REST, integridad y estructura de la base de datos, e importación de mazos desde MarvelCDB. Las pruebas del sistema de entrenamiento del modelo SVM y del despliegue en producción se describen en las secciones~\ref{sec:svm} y~\ref{sec:despliegue} respectivamente.

\subsection{Metodología de Pruebas}

Se ha adoptado un enfoque de \textbf{pruebas manuales funcionales continuas}. Este método se caracteriza por un ciclo iterativo de desarrollo y verificación:

\begin{enumerate}
    \item \textbf{Desarrollo iterativo:} Cada funcionalidad implementada se sometió a pruebas inmediatas para identificar y corregir errores en etapas tempranas.
    \item \textbf{Pruebas en tiempo real:} Se utilizó el servidor de desarrollo de Vite con \textit{Hot Module Replacement} (HMR) para observar y validar los cambios instantáneamente.
    \item \textbf{Cobertura de escenarios:} Se verificó cada funcionalidad bajo múltiples condiciones, incluyendo datos válidos, formatos incorrectos y casos límite.
    \item \textbf{Pruebas de integración:} Se comprobó la correcta comunicación del cliente con el backend propio y la API externa de MarvelCDB.
\end{enumerate}

Las herramientas principales utilizadas durante este proceso fueron las \textit{Chrome DevTools} para la depuración de red, el compilador de \textbf{TypeScript} para la detección estática de errores, \textbf{ESLint} para el mantenimiento de la calidad del código y la plataforma \textbf{Vercel} para las pruebas en entorno de producción.

\subsection{Pruebas Críticas y Resolución de Errores}

A continuación se describen las pruebas realizadas sobre las funcionalidades más complejas del sistema. En cada apartado se indica el objetivo de la prueba, los problemas detectados y las soluciones aplicadas.

\subsubsection{1. Autenticación y Autorización}

\textbf{Pruebas realizadas antes de producción y soluciones aplicadas:}

\begin{itemize}
    \item \textbf{Problema:} El cierre de sesión (\textit{logout}) fallaba en páginas protegidas específicas como el Perfil de Usuario.
    \item \textbf{Solución:} Se estandarizó la llamada a la función de logout incluyendo los parámetros de retorno requeridos (\texttt{logoutParams: \{ returnTo: window.location.origin \}}) en todos los componentes.

    \item \textbf{Problema:} Al acceder sin estar logueado a una ruta protegida (p. ej. ``Mis mazos''), la aplicación no redirigía de forma fiable al login y en algunos casos mostraba contenido vacío o errores.
    \item \textbf{Solución:} Se revisó el guard o el componente que protege las rutas para que compruebe el estado de autenticación antes de renderizar y redirija siempre a la URL de login de Auth0 cuando no hay sesión, usando el origen dinámico (\texttt{window.location.origin}).
\end{itemize}

\subsubsection{2. Validación de Datos}

Se verificó la integridad de los datos introducidos por los usuarios, probando límites de cartas y unicidad de nombres de mazos.

\textbf{Errores encontrados y solucionados:}

\begin{itemize}
    \item \textbf{Problema:} El sistema permitía la creación de mazos con nombres duplicados si diferían solo en mayúsculas o minúsculas.
    \item \textbf{Solución:} Se implementó una validación global que normaliza los nombres antes de la comparación, impidiendo duplicados semánticos en la base de datos.
\end{itemize}

\subsubsection{3. Navegación SPA e Integración}

Se comprobó la fluidez de la navegación sin recargas y la correcta comunicación con el servidor.

\textbf{Errores encontrados y solucionados:}

\begin{itemize}
    \item \textbf{Problema:} El uso de \texttt{window.location.href} provocaba recargas completas de la página, perdiendo el estado de la aplicación.
    \item \textbf{Solución:} Se sustituyeron todas las navegaciones imperativas por el hook \texttt{useNavigate} de React Router DOM, preservando el estado de la SPA.

    \item \textbf{Problema:} La lista de favoritos mostraba al creador como ``anónimo'' en ciertas vistas.
    \item \textbf{Solución:} Se optimizó la consulta de carga de favoritos para incluir una llamada adicional al endpoint de detalles del mazo (\texttt{getDeckById}), recuperando correctamente el nombre del autor.
\end{itemize}

\subsection{Validación de Endpoints REST}

Se verificó que todos los endpoints de la API responden correctamente y devuelven los códigos HTTP y la estructura JSON esperados.

\textbf{Pruebas realizadas:} Comprobación exhaustiva con Postman (GET, POST, PUT, DELETE); validación de códigos 200, 201, 400, 401, 404 y 500; verificación de la estructura JSON de las respuestas y del manejo de errores con mensajes informativos.

\textbf{Resultado:} No se detectaron incidencias. Todos los endpoints funcionan correctamente y manejan adecuadamente los casos de éxito y error.

\subsection{Integridad de la Base de Datos y Estructura del Esquema}

Se garantizó la consistencia y coherencia de los datos mediante las restricciones de integridad referencial en SQLite y se validó que las tablas tengan la estructura necesaria (columnas presentes y tipadas correctamente), con capacidad del sistema para adaptarse a cambios de esquema sin intervención manual.

\subsubsection*{Pruebas realizadas}

\textbf{Integridad referencial.}
\begin{itemize}
    \item Verificación de claves foráneas entre \texttt{users}, \texttt{decks}, \texttt{game\_configurations}, \texttt{user\_favorites} y \texttt{deck\_comments}: se comprobó que las referencias (\texttt{user\_id}, \texttt{deck\_id}, \texttt{villain\_id}) apuntan a tablas y columnas existentes y que no se pueden insertar filas con IDs inexistentes cuando \texttt{PRAGMA foreign\_keys = ON} está activo (configurado en \texttt{db\_utils} en cada conexión).
    \item Validación del borrado en cascada: se eliminó un usuario de prueba y se comprobó que se borraron sus partidas en \texttt{game\_configurations} y sus favoritos en \texttt{user\_favorites}; se eliminó un mazo y se verificó que se eliminaron sus partidas asociadas, sus favoritos y sus comentarios en \texttt{deck\_comments}.
    \item Comprobación de ausencia de registros huérfanos: tras las pruebas de cascada, se ejecutaron consultas de conteo sobre las tablas relacionadas para confirmar que no quedan filas con \texttt{user\_id} o \texttt{deck\_id} referenciando registros ya eliminados.
\end{itemize}

\textbf{Estructura de columnas y migraciones.}
\begin{itemize}
    \item Verificación dinámica con \texttt{PRAGMA table\_info()} en \texttt{cards}, \texttt{decks}, \texttt{game\_configurations} y \texttt{user\_favorites}: se comprobó que los endpoints y las funciones de esquema (\texttt{ensure\_*}) consultan la estructura antes de leer o escribir, evitando asumir columnas que podrían no existir en bases antiguas.
    \item Migración automática de columnas en \texttt{cards}: se probó sobre una base sin las columnas \texttt{marvelcdb\_code}, \texttt{is\_unique}, \texttt{deck\_limit}, \texttt{traits} y \texttt{text}; tras arrancar la aplicación, \texttt{ensure\_cards\_columns()} añadió las columnas faltantes y se repitió la importación de cartas y las búsquedas por código sin errores.
    \item Migración automática de columnas en \texttt{decks}: se probó sobre una base sin \texttt{hero\_name}, \texttt{aspect}, \texttt{hero\_id} ni \texttt{user\_id}; \texttt{ensure\_decks\_columns()} añadió las columnas y se verificó la creación y listado de mazos con los nuevos campos.
    \item Estandarización de \texttt{user\_id} como INTEGER: en bases donde \texttt{game\_configurations} o \texttt{user\_favorites} tenían \texttt{user\_id} como TEXT (llegando a almacenar \texttt{auth0\_id}), las consultas de partidas y favoritos fallaban o no devolvían datos correctos. Se comprobó que la migración automática (tabla temporal, conversión de valores y reemplazo de tabla) convierte \texttt{user\_id} a INTEGER y resuelve \texttt{auth0\_id} a \texttt{users.id} cuando es necesario; tras la migración, las consultas por usuario funcionaron correctamente.
    \item Validación del uso de IDs numéricos: se comprobó que los endpoints de mazos, partidas y favoritos usan \texttt{deck\_id} e \texttt{user\_id} como enteros en las cláusulas \texttt{WHERE} y en las respuestas JSON, sin mezclar tipos que pudieran provocar fallos de clave foránea o comparaciones incorrectas.
\end{itemize}

\textbf{Ejemplo concreto (tabla \texttt{cards}).} En una base de datos creada con una versión antigua del esquema, la tabla \texttt{cards} no tenía las columnas \texttt{marvelcdb\_code} ni \texttt{is\_unique}. Al desplegar la versión que integra la importación desde MarvelCDB, las consultas que buscaban cartas por código (\texttt{GET /api/cards/marvelcdb-code/\{code\}}) y las inserciones que guardaban el código y el flag \texttt{is\_unique} fallaban. Se comprobó que, en arranque, \texttt{ensure\_cards\_columns()} obtiene la estructura con \texttt{PRAGMA table\_info(cards)}, detecta las columnas ausentes y ejecuta \texttt{ALTER TABLE cards ADD COLUMN marvelcdb\_code VARCHAR(20)} y \texttt{ALTER TABLE cards ADD COLUMN is\_unique BOOLEAN DEFAULT 0}, además de crear el índice único en \texttt{marvelcdb\_code}. Tras esa migración automática, la importación de cartas y las búsquedas por código funcionaron correctamente sin intervención manual sobre la base de datos.

\textbf{Resultado:} No se detectaron incidencias tras aplicar las correcciones. La integridad referencial funciona correctamente, el borrado en cascada limpia los datos relacionados y el esquema se adapta automáticamente a las columnas requeridas, manteniendo la consistencia necesaria para el funcionamiento de la aplicación.

\subsection{Importación de Mazos desde MarvelCDB}
\label{sec:import-marvelcdb}

Se verificó de forma unificada la capacidad del sistema para importar mazos desde la API pública de MarvelCDB: obtención de mazos populares o por URL, procesamiento del JSON en el cliente, resolución de cartas y del héroe contra la base de datos local, envío al backend y persistencia correcta en la tabla \texttt{decks}. Se prestó especial atención a la consistencia de identificadores en toda la cadena, ya que durante el desarrollo surgieron varias incidencias relacionadas con el uso de \texttt{user\_id}, \texttt{deck\_id} y \texttt{hero\_id}.

\subsubsection*{Pruebas realizadas}

\begin{itemize}
    \item \textbf{Backend: obtención y procesamiento.} Se comprobó \texttt{fetch\_popular\_decks()} (o el equivalente usado para mazos públicos), el \textit{parsing} de la estructura de mazos devuelta por MarvelCDB, la función \texttt{get\_card\_name\_by\_code()} (o búsqueda por \texttt{marvelcdb\_code}) para traducir códigos de carta a la base local, y el flujo de \texttt{import\_decks()} cuando la importación se hace desde el servidor. Se validó el manejo de errores ante fallos de conexión o datos inválidos.

    \item \textbf{Frontend: conversión y envío.} Se validó la importación desde la interfaz de usuario: introducción de URL o selección de mazo desde MarvelCDB, conversión de los datos al formato esperado por la API propia, detección de cartas faltantes (y su solicitud al endpoint de importación de cartas), y envío mediante \texttt{POST /api/decks} con \texttt{hero\_name}, \texttt{hero\_id}, \texttt{aspect} y lista de cartas correctamente formateada. Se comprobó que el mazo creado se muestra correctamente y que el \texttt{id} devuelto por el backend se utiliza de forma coherente en la aplicación (navegación, favoritos, historial de partidas).

    \item \textbf{Identificación unívoca de cartas.} En fases iniciales las cartas se buscaban por nombre, lo que generaba ambigüedades y emparejamientos incorrectos. Se verificó que tanto el backend como el frontend utilizan el código único de MarvelCDB (\texttt{marvelcdb\_code}) para buscar y asociar cartas, eliminando conflictos por nombres duplicados o similares.

    \item \textbf{Detección correcta del héroe.} Se detectó que en ciertos mazos el héroe se identificaba erróneamente (por ejemplo, \textit{Captain Marvel} asignada como \textit{Adam Warlock} por coincidencia de códigos de investigador en la API de origen). Se comprobó que el mapeo prioriza los campos \texttt{hero\_code} y \texttt{hero\_name} de la respuesta de MarvelCDB y se estableció un sistema de respaldo (\textit{fallback}) robusto, de modo que el \texttt{hero\_name} y, cuando corresponde, el \texttt{hero\_id} asociados al mazo en la base local son correctos.

    \item \textbf{Aspecto del mazo.} Se validó que el aspecto no se importa únicamente desde el valor devuelto por MarvelCDB, sino que se verifica contra la base de datos local (por ejemplo, consultando las cartas del mazo y su facción/aspecto) para asignar el aspecto según la taxonomía interna y evitar incoherencias con las cartas almacenadas.

    \item \textbf{Consistencia de identificadores (\texttt{user\_id}, \texttt{deck\_id}, \texttt{hero\_id}).} Este aspecto requirió un trabajo considerable durante el desarrollo. Se realizaron las siguientes comprobaciones:
    \begin{itemize}
        \item \textbf{\texttt{user\_id}:} En la tabla \texttt{decks}, \texttt{user\_id} debe ser de tipo \texttt{INTEGER} (referencia a \texttt{users.id}). En versiones antiguas, \texttt{game\_configurations} y \texttt{user\_favorites} tenían \texttt{user\_id} como \texttt{TEXT}, llegando a almacenar el \texttt{auth0\_id} en lugar del ID numérico del usuario, lo que rompía las claves foráneas y las consultas que asocian mazos, partidas y favoritos al usuario. Se validó la migración automática de \texttt{user\_id} de \texttt{TEXT} a \texttt{INTEGER} en \texttt{game\_configurations} y \texttt{user\_favorites} (resolviendo \texttt{auth0\_id} a \texttt{users.id} cuando era necesario) y que \texttt{decks.user\_id} se cree siempre como \texttt{INTEGER}. En la importación, si el usuario no está autenticado, \texttt{user\_id} se persiste como \texttt{NULL} sin provocar errores.
        \item \textbf{\texttt{deck\_id}:} El backend devuelve el \texttt{id} del mazo recién creado como entero (\texttt{cursor.lastrowid}). Se comprobó que el frontend y el resto de endpoints (favoritos, partidas, comentarios) tratan \texttt{deck\_id} como \texttt{INTEGER} y que no se producen fallos por comparación o serialización como cadena.
        \item \textbf{\texttt{hero\_id}:} Es un campo opcional en \texttt{decks} que referencia al héroe en la base local. Se verificó que el frontend puede enviar \texttt{hero\_id} tras resolver el héroe por \texttt{hero\_code} o \texttt{hero\_name}, y que el backend lo persiste correctamente; si no se envía, el campo queda \texttt{NULL} y el mazo sigue siendo válido.
    \end{itemize}
\end{itemize}

\textbf{Resultado:} Tras aplicar las correcciones anteriores, la importación de mazos desde MarvelCDB funciona correctamente en frontend y backend. Los mazos se obtienen, se procesan, se crean en la base de datos con el dueño y el héroe correctos, y los identificadores se mantienen coherentes en todas las tablas y en la comunicación con el cliente, permitiendo poblar la base de datos con contenido de la comunidad y garantizando la integridad referencial y el correcto funcionamiento de favoritos, partidas y recomendaciones.

\subsection{Resultados y Lecciones Aprendidas}

El proceso de pruebas permitió identificar y corregir 12 errores críticos antes del lanzamiento, garantizando una aplicación estable y funcional.

Entre las lecciones aprendidas destacan:
\begin{itemize}
    \item La importancia de las \textbf{pruebas continuas} inmediatas tras el desarrollo para evitar la acumulación de deuda técnica.
    \item La utilidad del \textbf{tipado estático} de TypeScript para prevenir errores en tiempo de ejecución.
    \item La necesidad de validar siempre en el \textbf{entorno de producción}, ya que pueden surgir discrepancias de configuración (especialmente en redirecciones y variables de entorno) respecto al entorno local.
\end{itemize}

Como trabajo futuro, se recomienda la implementación de tests automatizados (unitarios y \textit{end-to-end}) para aumentar la cobertura y facilitar el mantenimiento a largo plazo.

% -----------------------------------------------------------------------
\section{Sistema de Entrenamiento Automático del Modelo SVM}
\label{sec:svm}
% -----------------------------------------------------------------------

Se validó que el modelo de IA se reentrena automáticamente tras registrar nuevas partidas, se persiste en disco y se carga correctamente en producción para generar recomendaciones. Durante las pruebas se identificaron y resolvieron varias incidencias relacionadas con la disponibilidad de las tablas de datos, la cantidad mínima de partidas, la persistencia del modelo en el entorno de producción y la localización del archivo del modelo.

\subsubsection*{Pruebas realizadas}

\begin{itemize}
    \item \textbf{Disparo del entrenamiento en segundo plano:} Se comprobó que, al registrar una nueva partida mediante \texttt{POST /api/game-configurations}, se encola una tarea en segundo plano (\texttt{BackgroundTasks}) que ejecuta el entrenamiento del SVM. De este modo la respuesta al usuario no se bloquea y el modelo se actualiza de forma asíncrona.

    \item \textbf{Existencia de tablas antes del entrenamiento:} En fases iniciales se produjeron fallos por ausencia de la tabla \texttt{game\_configurations} o de columnas necesarias en \texttt{decks} cuando el entrenamiento se ejecutaba nada más arrancar el servicio o en entornos recién desplegados. Se validó que, antes de leer datos para entrenar, se garantiza la existencia de las tablas (\texttt{ensure\_game\_configurations\_table}, \texttt{ensure\_decks\_columns}) y se verifica la presencia de \texttt{game\_configurations} y \texttt{decks} (\texttt{verify\_tables\_exist}) en el módulo de preparación de datos. Con estas comprobaciones, el entrenamiento no se ejecuta sobre un esquema incompleto y se evitan errores en tiempo de ejecución.

    \item \textbf{Cantidad mínima de datos y división entrenamiento/prueba:} Se verificó el comportamiento con pocas partidas. El entrenamiento requiere al menos dos partidas para poder aplicar una división 80/20 entre conjunto de entrenamiento y de prueba. Con una sola partida no se entrena y no se muestra error al usuario (comportamiento silencioso). Se comprobó además que, cuando hay partidas de ambos resultados (victorias y derrotas), se utiliza \texttt{stratify} en la división para mantener la proporción de clases; cuando solo hay un resultado (por ejemplo solo victorias), se entrena sin \texttt{stratify} para no provocar fallos. El cálculo de precisión en entrenamiento y en prueba se validó con la división 80/20.

    \item \textbf{Persistencia del modelo en producción:} Se validó que el modelo se guarda en un archivo \texttt{.pkl} en la ruta indicada por la variable de entorno \texttt{MODEL\_PATH}. En Render se configuró un volumen persistente montado en \texttt{/data} y \texttt{MODEL\_PATH} apuntando a un archivo dentro de ese volumen (por ejemplo \texttt{/data/saved\_models/villain\_svm\_model.pkl}), de modo que el modelo sobrevive a reinicios del servicio. Se comprobó mediante reinicios que el archivo \texttt{.pkl} permanece y que la aplicación carga el modelo correctamente para las recomendaciones.

    \item \textbf{Carga del modelo en tiempo real:} Se verificó que, ante una petición de recomendación (\texttt{POST /api/recommendations/deck}), el backend carga el modelo desde la ruta configurada por \texttt{MODEL\_PATH} cuando existe, y que en caso de no existir modelo (por ejemplo por datos insuficientes) se utiliza el \textit{fallback} determinista sin generar error al usuario.

    \item \textbf{Eliminación y localización del modelo:} En las pruebas de limpieza de datos de producción (\textit{reset}) se detectaron incidencias al intentar borrar el modelo: el script buscaba el archivo en rutas fijas que no coincidían con el entorno de Render o con el nombre configurado. Se validó que la localización del modelo se realiza de forma unificada mediante \texttt{MODEL\_PATH} en todo el código (entrenamiento, carga para recomendaciones y script de reset), de modo que no queden rutas \textit{hardcodeadas} y el comportamiento sea correcto en desarrollo y en producción.

    \item \textbf{Consistencia de tipos y tablas:} Se comprobó que el entrenamiento utiliza correctamente \texttt{user\_id} y \texttt{deck\_id} como enteros (\texttt{INTEGER}) en \texttt{game\_configurations} y \texttt{decks}, en línea con el diseño de la base de datos, evitando inconsistencias que en el pasado afectaron a consultas relacionadas con el historial de partidas.
\end{itemize}

\textbf{Resultado:} Tras aplicar las correcciones anteriores, no se detectaron incidencias en el funcionamiento del sistema de entrenamiento automático. El modelo se actualiza automáticamente con nuevos datos, se persiste en disco en la ruta configurada y mantiene su estado entre reinicios y despliegues en Render. Las recomendaciones utilizan el modelo cuando está disponible y recurren al \textit{fallback} determinista cuando no hay modelo entrenado, sin afectar la experiencia del usuario.

% -----------------------------------------------------------------------
\section{Pruebas de Despliegue en Producción}
\label{sec:despliegue}
% -----------------------------------------------------------------------

Se verificó el correcto funcionamiento del backend en el entorno de nube (Render.com): persistencia de la base de datos SQLite y del modelo SVM entre reinicios, carga correcta de la configuración y comunicación con el frontend desplegado en Vercel.

\textbf{Pruebas realizadas:} Configuración y montaje del volumen persistente en \texttt{/data}; lectura de las variables de entorno \texttt{DB\_PATH} y \texttt{MODEL\_PATH}; inicialización automática de tablas en arranque en frío; importación masiva inicial de cartas (1955+ registros); configuración CORS para el dominio del frontend; comprobación tras reinicio del servicio de que la base de datos y el modelo siguen disponibles.

\textbf{Errores encontrados y solucionados:}

\begin{itemize}
    \item \textbf{Problema:} Tras un reinicio del servicio, la base de datos y el modelo desaparecían porque se usaba el sistema de archivos efímero del contenedor.
    \item \textbf{Solución:} Se configuró un disco persistente en Render montado en \texttt{/data} y se definieron \texttt{DB\_PATH} y \texttt{MODEL\_PATH} apuntando a rutas dentro de ese directorio, de modo que los datos sobreviven a los despliegues y reinicios.

    \item \textbf{Problema:} Las peticiones desde el frontend (Vercel) al backend (Render) fallaban por políticas CORS (origen no permitido).
    \item \textbf{Solución:} Se añadieron los orígenes del frontend en producción a la lista de permitidos del middleware CORS de FastAPI, incluyendo el dominio de Vercel y, si aplica, variantes con y sin \texttt{www}.

    \item \textbf{Problema:} En arranque en frío, la aplicación intentaba acceder a la base de datos o importar cartas antes de que el volumen estuviera disponible, generando errores de archivo no encontrado.
    \item \textbf{Solución:} Se aseguró que la inicialización de tablas y la importación de cartas se ejecutan tras comprobar la existencia del directorio \texttt{/data} (o de las rutas configuradas) y se manejaron las excepciones para reintentar o informar con claridad en los logs.

    \item \textbf{Problema (URLs):} En producción, el frontend realizaba las llamadas a la API usando una URL incorrecta (por ejemplo \texttt{localhost} o una URL antigua del backend), de modo que las peticiones fallaban por red o devolvían errores 404. Al cambiar la URL del servicio en Render o al desplegar por primera vez, la aplicación cliente no apuntaba al backend correcto.
    \item \textbf{Solución:} Se definió la URL base del backend como variable de entorno en el frontend (por ejemplo \texttt{VITE\_API\_URL} en Vite) y se configuró en Vercel con la URL real del servicio en Render (p. ej. \texttt{https://marvelcdb-backend.onrender.com}). Así, el frontend usa en producción la URL correcta del API sin hardcodearla y se evitan fallos al cambiar de entorno o al redeployar el backend.

    \item \textbf{Problema (autenticación):} Tras el despliegue, el inicio de sesión y el registro fallaban con error \textit{Callback URL mismatch}: Auth0 rechazaba la redirección porque la URL de retorno configurada en el panel no coincidía con el dominio real del frontend en Vercel.
    \item \textbf{Solución:} Se añadieron en el panel de Auth0 las URLs de producción (dominio de Vercel) en \textit{Allowed Callback URLs}, \textit{Allowed Logout URLs} y \textit{Allowed Web Origins}, y se aseguró que el código use el origen dinámico (\texttt{window.location.origin}) para las redirecciones, de modo que funcione en cualquier dominio desplegado.
\end{itemize}

\textbf{Resultado:} Tras aplicar las correcciones, el despliegue en producción es estable; la base de datos y el modelo persisten correctamente, el frontend se comunica con el backend sin errores CORS ni por URL incorrecta, y el inicio de sesión y registro funcionan correctamente en el dominio de producción.
